\begin{figure}[ht]
\centering
%%%%%%%%%%%%%%%%%%%%%%%%%%%%%%%%%%%%%%%% left: profile along the cut A--A'
\begin{tikzpicture}[x=1.15cm,y=1.65cm,>={Stealth[length=5pt]},
  declare function={
    plo(\x) = 0.85 + 0.75*exp(-(\x-1.2)^2/0.45) + 0.55*exp(-(\x-2.9)^2/0.5);
    pup(\x) = 2.55 - 0.90*exp(-(\x-4.6)^2/0.55) - 0.55*exp(-(\x-2.9)^2/0.45);
  }]
  % the cut A--A', carrying the traces of the contact sets
  \draw[gray] (0,0.50) -- (6,0.50);
  \node[gray,left] at (0,0.50) {$A$};
  \node[gray,right] at (6,0.50) {$A'$};
  \draw[Blue,line width=2pt] (0.75,0.50) -- (1.7,0.50);
  \draw[Red,line width=2pt] (4.2,0.50) -- (5.0,0.50);
  \node[Blue,below] at (1.225,0.48) {$\Cbot$};
  \node[Red,below] at (4.6,0.48) {$\Ctop$};
  \draw[dotted,gray] (0.75,0.50) -- (0.75,{plo(0.75)});
  \draw[dotted,gray] (1.7,0.50) -- (1.7,{plo(1.7)});
  \draw[dotted,gray] (4.2,0.50) -- (4.2,{pup(4.2)});
  \draw[dotted,gray] (5.0,0.50) -- (5.0,{pup(5.0)});
  % obstacles
  \draw[Blue,thick] plot[domain=0:6,samples=200,smooth] (\x,{plo(\x)});
  \draw[Red,thick]  plot[domain=0:6,samples=200,smooth] (\x,{pup(\x)});
  \node[Blue,below] at (3.3,{plo(3.3)}) {$\psilo$};
  \node[Red,above] at (1.1,{pup(1.1)}) {$\psiup$};
  % solution, heavy where it coincides with an obstacle
  \draw[line width=1.2pt] (0,1.95) to[out=-58,in=230] (0.75,{plo(0.75)});
  \draw[line width=2pt] plot[domain=0.75:1.7,samples=60,smooth] (\x,{plo(\x)});
  \draw[line width=1.2pt] (1.7,{plo(1.7)}) to[out=-41,in=196] (2.15,1.31)
                                           to[out=16,in=205] (2.9,1.53)
                                           to[out=25,in=176] (3.70,1.95)
                                           to[out=-4,in=132] (4.2,{pup(4.2)});
  \draw[line width=2pt] plot[domain=4.2:5.0,samples=60,smooth] (\x,{pup(\x)});
  \draw[line width=1.2pt] (5.0,{pup(5.0)}) to[out=50,in=192] (6,2.22);
  \node[above] at (0.30,1.90) {$u$};
  % the uniform gap
  \draw[<->,thick] (2.9,{plo(2.9)}) -- (2.9,{pup(2.9)});
  \node[left,fill=white,inner sep=1pt] at (2.9,1.80) {$\delta_{\min}$};
\end{tikzpicture}
\hfill
%%%%%%%%%%%%%%%%%%%%%%%%%%%%%%%%%%%%%%%% right: plan view of Omega
\begin{tikzpicture}[x=1.42cm,y=1.42cm,>={Stealth[length=5pt]}]
  \draw[thick] (0,0) rectangle (4.6,3.2);
  \node at (0.35,2.9) {$\Om$};
  % the cut A--A'
  \draw[gray] (0,2.05) -- (4.6,2.05);
  \node[gray,left] at (0,2.05) {$A$};
  \node[gray,right] at (4.6,2.05) {$A'$};
  % contact sets
  \filldraw[Blue,fill=Blue!12,thick]
    plot[smooth cycle] coordinates {(0.55,1.55) (0.75,2.25) (1.45,2.45) (1.95,1.85) (1.75,1.05) (1.0,0.95)};
  \filldraw[Red,fill=Red!12,thick]
    plot[smooth cycle] coordinates {(2.95,1.35) (3.1,2.05) (3.8,2.35) (4.2,1.75) (3.95,1.05) (3.25,0.9)};
  \draw[Blue,line width=2pt] (0.65,2.05) -- (1.85,2.05);
  \draw[Red,line width=2pt] (3.12,2.05) -- (4.0,2.05);
  \node[Blue] at (1.2,1.6) {$\Cbot$};
  \node[Red] at (3.6,1.62) {$\Ctop$};
  % separation
  \draw[<->,thick] (1.9,1.15) -- (2.98,1.15);
  \node[below] at (2.44,1.12) {$\ge r_0$};
\end{tikzpicture}
\caption{A notional illustration of Proposition~\ref{prop:separation}.  Left: the profile along the cut $A$--$A'$, with the solution $u$ drawn heavy where it coincides with an obstacle.  Right: a plan view of $\Om$.  The Proposition converts the vertical gap $\delta_{\min}$ into the horizontal separation $r_0$.}
\label{fig:separation}
\end{figure}
